\documentclass[12pt]{article}

\usepackage{sbc-template}
\usepackage{graphicx,url}
\usepackage[utf8]{inputenc}
\usepackage[brazil]{babel}

\sloppy

\title{Mitigação de Viés Algorítmico em Cenários de Desbalanceamento Sociodemográfico: Uma Análise Comparativa entre Foundation Models (TabPFN) e Gradient Boosting (XGBoost) nas Eleições Municipais do Ceará\\ Papers and Abstracts}

\author{Luana Teles Alves\inst{1}, Luann Alves Pereira de Lima\inst{1}, Thiago Nerton Macedo\\ 
  Alves\inst{1}, Vitória do Nascimento Pontes\inst{1} }

\address{Centro de Ciências e Tecnologia (CCT) -- Universidade Federal do Cariri (UFCA)\\
  Juazeiro do Norte -- CE -- Brasil
  \email{luana.teles@aluno.ufca.edu.br, luann.alves@aluno.ufca.edu.br,\\
  thiago.nerton@aluno.ufca.edu.br, vitoria.pontes@aluno.ufca.edu.br}
}

\begin{document} 

\maketitle

\begin{abstract}
  This meta-paper describes the style to be used in articles and short papers
  for SBC conferences. For papers in English, you should add just an abstract
  while for the papers in Portuguese, we also ask for an abstract in
  Portuguese (``resumo''). In both cases, abstracts should not have more than
  10 lines and must be in the first page of the paper.
\end{abstract}
     
\begin{resumo} 
  Sob validação cruzada ($K=10$, $N=2.178$) com atributos demográficos, ambos
  atingiram desempenho preditivo equivalente (MCC $\approx$ 0,28). Contudo, o
  XGBoost apresentou viés de gênero e raça significativo (DIR = 1,17 e 0,88,
  respectivamente), amplificando padrões históricos de auto-seleção. Em
  contrapartida, o TabPFN manteve paridade relativa em ambas as dimensões
  (DIR $\approx$ 1,01 e 0,91), evidenciando que diferenças arquiteturais entre
  foundation models e gradient boosting podem gerar disparidades de equidade
  relevantes mesmo sob acurácia preditiva equivalente.
\end{resumo}


\section{Introdução}

A predição de resultados eleitorais em eleições majoritárias no Ceará enfrenta desafios
metodológicos complexos dada a natureza desequilibrada dos dados – muitos candidatos para
poucas vagas – e a influência de barreiras históricas sociodemográficas. O cargo de
prefeito, em particular, exige uma análise que vá além da simples performance estatística,
visando na equidade e integridade dos modelos de decisão.

Com o  uso dos algoritmos de aprendizado de máquina cada vez mais presente nos processos de 
tomadas de decisões em diferentes domínios sociais, surge a preocupação urgente em relação a
amplificação de preconceitos contra um grupo de pessoas, baseados em características sensíveis
como raça e gênero. Modelos baseados em árvores de decisão, frequentemente utilizados em
dados tabulares, operam sob a lógica de otimização da acurácia global, o que pode resultar
em padrões de classificação distintos para diferentes grupos populacionais \cite{grari:2020}. 

Historicamente algoritmos que se baseiam no Gradient Boosting, como o XGBoost, viraram padrão
na indústria quando se trata de dados tabulares devido a suas taxas de acurácia geralmente mais
altas. Sua eficácia reside na construção iterativa de árvores de decisão, que, através da
otimização de uma função de perda global, conseguem capturar relações não-lineares complexas
com alta eficiência computacional \cite{ravichandran:2020}. No entanto, esse paradigma de "minimização de perda empírica"
impõe limitações estruturais: modelos dependentes exclusivamente de conjuntos de dados locais são
suscetíveis à cristalização de vieses presentes no histórico de treinamento, além de exigirem uma
otimização exaustiva de hiperparâmetros para alcançar estabilidade preditiva em “small data”.

Nesse cenário, emerge o paradigma dos Foundation Models para dados tabulares, representados de forma
proeminente pelo TabPFN (Prior-Data Fitted Network). Diferente dos métodos de boosting, o TabPFN é
uma rede neural baseada na arquitetura Transformer que realiza inferência através de Aprendizado em
Contexto (In-Context Learning) \cite{garg:2025,qu:2025}. Em vez de aprender os pesos durante um processo de treinamento
interativo, o modelo utiliza um conhecimento prévio consolidado durante um vasto treinamento em datasets
sintéticos e reais, permitindo que a predição em novos dados seja executada como uma consulta bayesiana
sobre as distribuições previamente aprendidas [4, 9]. Embora a literatura internacional apresenta avanços
significativos na mitigação desses vieses em aprendizado de máquina, sua aplicação prática na área eleitoral
em específico é baixa ou quase nula. Estudos recentes têm comparado o desempenho de arquiteturas de fundação
e métodos tradicionais em benchmarks sintéticos \cite{robertson:2024,garg:2025}, porém, carecem de validação em contextos reais de
eleições municipais brasileiras. 

Neste artigo, o objetivo central é realizar uma análise comparativa do desempenho preditivo e das métricas
de equidade entre o TabPFN e o XGBoost, utilizando dados do Tribunal Superior
Eleitoral \cite{basedosdados:2026}. A investigação baseia-se em uma abordagem de validação cruzada estratificada
(StratifiedKFold, $K=10$), buscando observar como essas arquiteturas se comportam frente ao desbalanceamento
das classes e à presença de atributos sensíveis.  Com base nas propriedades arquiteturais contrastantes
entre o TabPFN e o XGBoost, formulam-se as seguintes hipóteses de trabalho: H1 (Paridade de Performance):
Ambas as arquiteturas apresentam desempenho preditivo estatisticamente equivalente em cenários de small
data eleitoral, conforme mensurado pelo Coeficiente de Correlação de Matthews (MCC). H2 (Divergência de
Equidade de Gênero): O XGBoost apresenta viés de gênero mensurável ($DIR \neq 1$), enquanto o TabPFN,
devido ao seu paradigma de inferência bayesiana amortizada, não evidencia disparidade estatisticamente
significativa nesta dimensão. H3 (Divergência de Equidade Racial): O XGBoost amplifica desigualdades
raciais latentes nos dados históricos, enquanto o TabPFN atenua tais disparidades, mantendo os indicadores
de equidade próximos à neutralidade ($DIR \approx 1$). Os resultados obtidos indicam que, embora ambas as
arquiteturas alcancem desempenho preditivo estatisticamente equivalente, o XGBoost apresenta viés de
gênero e raça confirmados por indicadores de Impacto Disparado (DIR). Em contraste, o TabPFN não evidencia
disparidades significativas, sugerindo que as diferenças arquiteturais entre modelos de fundação e métodos
de Gradient Boosting podem produzir impactos distintos na equidade algorítmica, mesmo sob acurácia preditiva
equiparável O restante deste artigo está estruturado da seguinte forma: a Seção 3 detalha o protocolo
experimental, incluindo a descrição dos dados do TSE e o desenho do split temporal. A Seção 4 apresenta
os resultados quantitativos, com foco nas métricas de performance e justiça algorítmica. Por fim, a Seção
5  discute as implicações desses achados para a integridade do processo eleitoral, apresentando as
considerações finais e perspectivas para trabalhos futuros.

\section{Metodologia} \label{sec:metodologia}

\subsection{Fonte de dados}

Os dados utilizados neste estudo foram obtidos diretamente do repositório de dados
abertos do Tribunal Superior Eleitoral (TSE) \cite{basedosdados:2026}, órgão responsável pela organização
e fiscalização do processo eleitoral brasileiro. Para fins desta pesquisa, foi feito
um recorte temporal compreende quatro ciclos eleitorais municipais, entre 2012 e 2024,
e aplicou-se um filtro geográfico restrito ao estado do Ceará. Além disso, foram mantidas
apenas as candidaturas ao cargo de prefeito municipal, o que delimita um contexto de small
data, no qual o modelo TabPFN costuma se destacar em relação ao XGBoost.

\subsection{Unidade de Análise}
A unidade de análise adotada é o indivíduo candidato, de modo que cada instância do conjunto
de dados corresponde a um registro único de candidatura. O conjunto de dados totaliza 2.216
instâncias, distribuídas ao longo dos quatro pleitos analisados. Para cada candidato, o
conjunto de dados contém informações referentes ao município de disputa, características
sociodemográficas (cor/raça autodeclarada, estado civil, faixa etária, gênero e grau de
instrução), dados político-partidários (partido de filiação e número de candidatura),
ocupação declarada, situação final na totalização dos votos, quantitativo de votos válidos,
votos nominais recebidos e data de carga do registro no sistema do TSE.

\subsection{Variáveis}

As variáveis foram organizadas em dois grupos funcionais. O primeiro reúne as características
sociodemográficas e político-institucionais do candidato que serão utilizadas como preditores:
gênero, raça/cor autodeclarada, faixa etária, estado civil, grau de instrução, ocupação declarada,
partido político de filiação, número do candidato e município. O segundo grupo diz respeito à
variável dependente, a situação da totalização indicando se aquele indivíduo foi eleito ou não.

Variáveis que introduzem data leakage foram excluídas do processo de modelagem. Em particular,
os campos de votos válidos e votos nominais foram descartados, uma vez que constituem informações
geradas pelo próprio processo eleitoral que se busca predizer. Afinal, sua inclusão como preditores
produziria modelos trivialmente precisos, porém sem qualquer capacidade de generalização para
cenários de predição pré-eleitoral \cite{kaufman:2012}.

\subsection{Distribuições relevantes}

O conjunto de dados apresenta desbalanceamento de classes expressivo na variável-alvo: dos 2.216
registros, 1.411 correspondem a candidatos não eleitos (63,7%) e 767 a candidatos eleitos (34,6%)
— o restante corresponde aos candidatos em segundo turno que são removidos durante a etapa de
pré-processamento. Nesse sentido, a disparidade da variável é esperada, tendo em vista a própria
estrutura competitiva das eleições municipais, nas quais cada disputa comporta apenas um vencedor
por município. Esse desequilíbrio foi considerado nas etapas de treinamento e avaliação dos modelos,
conforme descrito nas subseções subsequentes.

No que tange à composição de gênero, observa-se que 1.883 candidatos são do sexo masculino (85%)
e 333 do sexo feminino (15%), distribuição que reflete a histórica sub-representação feminina nos
cargos executivos municipais brasileiros e que se mantém consistente com padrões identificados em
estudos anteriores sobre o perfil das candidaturas no país [12]. Entre as candidatas do sexo feminino,
a taxa de êxito eleitoral é proporcionalmente inferior à dos candidatos masculinos, o que pode levar
os modelos a carregar esse preconceito em suas predições \cite{grari:2020}.

\section{Referências Bibliográficas}

Bibliographic references must be unambiguous and uniform.  We recommend giving
the author names references in brackets, e.g. \cite{knuth:84},
\cite{boulic:91}, and \cite{smith:99}.

The references must be listed using 12 point font size, with 6 points of space
before each reference. The first line of each reference should not be
indented, while the subsequent should be indented by 0.5 cm.

\bibliographystyle{sbc}
\bibliography{sbc-template}

\end{document}
